% !TEX root = ../Thesis.tex
\chapter{TESTING AND RESULTS}
\section{Current verification status}
A production build of the React application was successfully generated with Vite on 17 October 2026. This verifies that the current build command completed, but it does not establish usability, accessibility, or educational effectiveness.

\begin{table}[htbp]
\centering
\caption{Current verification status and planned evaluation activities.}
\label{tab:verification}
\begin{tabularx}{\textwidth}{|X|p{3cm}|X|}
\hline
\textbf{Activity} & \textbf{Status} & \textbf{Evidence / required action}\\ \hline
Production build & Passed & Vite production build completed successfully on 17 October 2026.\\ \hline
Keyboard accessibility & Pending & Test focus order and all interactive controls.\\ \hline
Screen-reader usability & Pending & Record screen-reader and browser setup plus observations.\\ \hline
Participant usability test & Pending & Recruit participants and record task outcomes.\\ \hline
Learning effectiveness & Pending & Collect and analyze pre/post knowledge responses.\\ \hline
\end{tabularx}
\end{table}

\section{Usability evaluation results}
\textit{To be completed after participant testing.} Record the methods, number of participants, tasks, completion rates, observed issues, and participant feedback. Do not replace missing results with estimates.

\section{Accessibility evaluation results}
\textit{To be completed after accessibility testing.} Document tested WCAG criteria, devices, browsers, assistive technologies, findings, and outstanding issues.

\section{Learning evaluation results}
\textit{To be completed after data collection.} Present anonymized pre-test and post-test results, sample size, question design, and appropriate limitations. Conclusions about learning improvement must follow from the actual data.
